\section{conic と二直線の三重点：短い排除証明書}

対象は $C=A_{012}\cup L_{03}\cup L_{04}$。
conic-plane 補題により、conic の bracket matrix の内部 normal 成分
$B_{00}=-z_{31}$ は非零の定数である。特に三重点の $\pi_1$ は非零。
三重点の slope 消滅証明の $\pi_1\ne0$ の部分を適用できる。
その証明で $k=1$ を使うのは conormal bracket が交点で零となる枝だけであるが、
conic の bracket は非零なので、それらは常に射影直線である。
従って三つの modular slopes は全て零、すなわち
\[
z_{37}=z_{39}=z_{41}=0.
\]
各 residue は active なので、その二次係数は非零である。
Poisson 全体の非零定数倍と $\operatorname{diag}(1,r,r^2,s,t)$ の射影変換により
\[
z_{38}=z_{40}=z_{42}=z_{31}=1
\]
と正規化できる。具体的に初期の二次係数を $\beta_c,\beta_1,\beta_2$ とすれば、
Poisson 倍率 $c$ に対して $r=c\beta_c,s=c\beta_1,t=c\beta_2$ を選ぶ。
そのとき $z_{31}'=c\beta_1\beta_2 z_{31}/\beta_c^2$ なので、四条件を同時に満たせる。

\subsection{rank-one 条件だけによる初期消去}

conic の bracket matrix の第一列は
\[
\begin{pmatrix}-1\\t^2z_{34}+tz_{32}\\t^2z_{35}+tz_{33}\end{pmatrix}.
\]
rank one から残る成分はこの列の積で決まり、係数比較により
\[
z_{23}=z_{26}=z_{27}=0,\quad
z_{24}=-z_{33}^2/3,\quad z_{28}=z_{32}^2/3,
\]
\[
z_{29}=z_{32}z_{33}/3-2z_{36}/3,
\]
\[
z_{13}=z_{34}^2,\quad z_{14}=z_{34}z_{35},\quad z_9=-z_{35}^2,
\]
\[
z_{11}=2z_{20}+2z_{32}z_{34},\quad
z_8=2z_{18}-2z_{33}z_{35},\quad
z_{12}=2z_{21}+z_{32}z_{35}+z_{33}z_{34}.
\]
この下で第一の直線の bracket matrix は
\[
B_1(t)=
\begin{pmatrix}
0&tz_{32}^2/3&tz_{30}\\
tz_{32}^2/3&-tz_{20}-z_{32}/2&-tz_{22}\\
tz_{30}&-tz_{22}&tz_{16}
\end{pmatrix}.
\]
第二の直線では $z_{32},z_{30},z_{20},z_{22}$ をそれぞれ
$z_{33},z_{25},z_{18},z_{19}$ に置き換えた形であり、必要な符号だけが変わる。
対称 rank-one 行列で第一対角成分が零なら、第一行・列は零である。
従って複素閉点では
\[
z_{32}=z_{30}=z_{33}=z_{25}=0.
\]
特に $z_{24}=z_{28}=0$、$z_{36}=-3z_{29}/2$ となる。

両直線の tangent Hamiltonian 行列の定数項は、normal 基底の第二方向が
conic 内部 normal $q$、第三方向が他方の直線の tangent となる順序で
\[
H_1(0)=-E_{23},\qquad H_2(0)=E_{23}.
\]
中心固有値条件 $(H_i+R_i I/2)B_i=0$ の $t$ 係数を比較すると
\[
z_{22}=z_{16}=z_{19}=0.
\]
従って
\[
B_1(t)=-tz_{20}E_{22},\qquad B_2(t)=-tz_{18}E_{22}.
\]
一般点で Heisenberg なので $z_{20},z_{18}$ は非零である。

\subsection{二つの Jacobi 係数と一つの中心条件}

第21版の全 Jacobi 方程式の第19番と第143番（零始まり）を上の必要条件で特殊化すると、
それぞれ
\[
J_{19}=z_{29}/5,\qquad J_{143}=z_5-z_{29}z_{34}
\]
となる。従って $z_{29}=z_5=0$。
一方、第一直線の中心条件の第二対角成分の $t^3$ 係数は
\[
E_3=z_{20}(z_5-3/10).
\]
これが零で、$z_{20}\ne0$ なので $z_5=3/10$。矛盾である。

全てを polynomial-combination 証明書として書けば
\[
E_3-z_{20}J_{143}-5z_{20}z_{34}J_{19}
=-\frac3{10}z_{20}.
\]
従って Jacobi と中心条件は $z_{20}=0$ を強制し、一般 Heisenberg 条件と両立しない。
この証明ではパラメータで割る必要はなく、rank-one 条件の閉点での帰結と、
明示した非零条件だけを用いている。
